% ============================================================================
% CONCLUSION DU CHAPITRE 4 
% source    : partie4/conclusion_partie3.tex
% ============================================================================

\section{CONCLUSION DU CHAPITRE~\ref{chap:resultats_discussions}}
\label{sec:conclusion_resultats_discussions}

\begin{sloppypar}
Ce chapitre a soumis le cadre th\'eorique
construit aux chapitres~\ref{chap:cadre_formel} et~\ref{chap:constructions} \`a l'\'epreuve de la coh\'erence
interne et de l'interpr\'etation. Quatre r\'esultats structurent le bilan.

Le premier est une calibration. Les ordres de grandeur retenus pour
les syst\`emes \'etudi\'es sont mutuellement coh\'erents et plausibles au regard
de la litt\'erature cognitive. Cette coh\'erence ne vaut pas validation, car les
valeurs ont \'et\'e fix\'ees par calibration r\'etrospective et les syst\`emes
d\'eploy\'es ne sont pas des instances du Trilemme
(Remarque~\ref{rem:portee_cas_etude}). De plus, les produits post-quantiques
illustratifs d\'epassent la majoration.

Le deuxi\`eme est \'epist\'emologique. La formalisation Lean4 et la relecture
des d\'efinitions ont conduit \`a pr\'eciser quatre points, \`a savoir la n\'ecessit\'e du r\'egime de
confidentialit\'e, le rapport entre la formulation initiale de la borne et sa
formulation certifi\'ee, la distinction entre fiabilit\'e de reconstruction et
fiabilit\'e de v\'erification ainsi que le caract\`ere non r\'ealisable du t\'emoin de
coh\'erence. Ces pr\'ecisions d\'elimitent la port\'ee des r\'esultats et montrent
l'int\'er\^et de combiner certification m\'ecanique et examen des d\'efinitions.

Le troisi\`eme est structurel. Sous ses hypoth\`eses, le Trilemme CRO
formalise la tension que le paradoxe de l'authenticit\'e invisible
d\'ecrivait. Une opposabilit\'e \'elev\'ee exige que la preuve informe le verdict,
ce qui, pour un interpr\`ete admissible, passe par de l'information sur le
message ou par la capacit\'e $\Delta$, alors que l'information sur le message d\'egrade la
confidentialit\'e. La forme raffin\'ee pr\'ecise quel facteur limite
l'opposabilit\'e selon le r\'egime. L'extension post-quantique montre que la
structure de la borne est la m\^eme pour des preuves quantiques. Elle montre aussi que le
terme de perte $\eta_q$ est nul pour les signatures post-quantiques
actuelles, qui sont des cha\^{\i}nes de bits.

Le quatri\`eme est prospectif. La cryptographie institutionnelle dispose d'un
formalisme, d'une borne prouv\'ee sous hypoth\`eses, de constructions et d'une
certification m\'ecanique, mais il lui manque des mesures. Quatre directions
restent ouvertes, en l'occurrence la calibration empirique de $\Delta$ et $\Delta_q$ par des
\'etudes contr\^ol\'ees, avec les crit\`eres de r\'efutation fix\'es au
chapitre~\ref{chap:experimentation}~; l'extension du cadre aux protocoles
interactifs et multi-parties~; la preuve ou la r\'efutation de la conjecture
d'optimalit\'e des constructions SOTP~; enfin la compl\'etion de la couverture
formelle Lean4. Ces ouvertures d\'elimitent le p\'erim\`etre de ce travail et
tracent le programme d'un sous-domaine dont le d\'eveloppement commence.

La conclusion g\'en\'erale dresse le bilan synth\'etique de l'ensemble.
\end{sloppypar}
